\subsection{算例 C3（$w=\SI{1}{\micro m}$，$L=\SI{15}{\micro m}$，区 II）的逐步代入}
\textbf{第 1 步：几何与孔体速率。}网格中喉道壁到壁长度 $L_m=15.031$~\si{\micro m}，喉口弧长 $2R\arcsin(h/R)$ 从孔壁中扣除：
\begin{align*}
A_S&=\pi\times5^2=78.540, & P_S&=2\pi\times5-1.0017=30.414,\\
A_B&=\pi\times20^2=1256.64, & P_B&=2\pi\times20-1.0001=124.664,\\
R_S&=\frac{0.005\times30.414}{78.540}+\frac{1}{3000}=2.270\times10^{-3}\ \mathrm{ms^{-1}}, & R_B&=8.294\times10^{-4}\ \mathrm{ms^{-1}}.
\end{align*}
\textbf{第 2 步：喉道。}$h=\SI{0.5}{\micro m}$：
\[
R_t=\frac{0.005}{0.5}+\frac{1}{3000}=0.01033,\qquad q=\sqrt{R_t/D}=\sqrt{0.01033/2.0}=0.07188\ \mathrm{\mu m^{-1}},\qquad qL=1.0804.
\]
\textbf{第 3 步：交换导纳与汇导纳（式 \eqref{eq:gx}）。}
\[
G_x=\frac{Dwq}{\sinh qL}=\frac{2.0\times1\times0.07188}{\sinh 1.0804}=0.1103,\qquad
G_s=Dwq\tanh\frac{qL}{2}=0.07090\ \mathrm{\mu m^2/ms}.
\]
\textbf{第 4 步：有效两点交换参数。}
\begin{align*}
R'_S&=R_S+\frac{G_s}{A_S}=2.270\times10^{-3}+9.027\times10^{-4}=3.172\times10^{-3}, & R'_B&=8.858\times10^{-4},\\
k_S&=\frac{G_x}{A_S}=1.404\times10^{-3}, & k_B&=\frac{G_x}{A_B}=8.778\times10^{-5},\\
\Psi_{\mathrm{eff}}&=\frac{k_S}{R'_S-R'_B}=\frac{1.404\times10^{-3}}{2.286\times10^{-3}}=0.6143. &&
\end{align*}
\textbf{第 5 步：本征速率（式 \eqref{eq:lam}）。}$a=R'_S+k_S=4.577\times10^{-3}$，$d=R'_B+k_B=9.736\times10^{-4}$：
\begin{align*}
\lambda_{f,s}&=\frac{a+d\pm\sqrt{(a-d)^2+4k_Sk_B}}{2}=\frac{5.550\times10^{-3}\pm3.671\times10^{-3}}{2}\\
&\Rightarrow\;T_{2,f}=1/\lambda_f=216.9\ \mathrm{ms},\qquad T_{2,s}=1/\lambda_s=1064.2\ \mathrm{ms}.
\end{align*}
\textbf{第 6 步：快模态振幅（式 \eqref{eq:amp}）。}
\[
\bar R'=p_SR'_S+p_BR'_B=1.020\times10^{-3},\qquad A_f=\frac{\bar R'-\lambda_s}{\lambda_f-\lambda_s}=0.0220,\qquad \eta=\frac{A_f}{p_S}=0.373.
\]
\textbf{第 7 步：对照。}TL 精确根给出 $T_{2,f}=228.8$~ms、$A_f=0.0225$；
NNLS（L-curve，SNR 1000）小峰 $T_{2\mathrm{lm}}=214.8$~ms、分数 0.0231；手算与精确根的差别来自准静态近似中 $q$ 取 $\lambda=0$。

\subsection{快交换算例 C5（$w=\SI{4}{\micro m}$，$L=\SI{2}{\micro m}$，区 III）的渐近式}
\[
\Psi_{\mathrm{eff}}=28.37,\qquad A_f\approx\frac{p_Sp_B^3}{\Psi_{\mathrm{eff}}^2}=\frac{0.0588\times0.9412^3}{28.37^2}=6.09\times10^{-5}\quad(\text{TL 精确}\ 5.50\times10^{-5}),
\]
\[
T_{2,f}\approx\frac{p_BA_S}{G_x}=\frac{0.9412\times78.54}{3.1728}=23.3\ \mathrm{ms}\quad(\text{TL}\ 23.0\ \mathrm{ms},\ \text{NNLS}\ 15.7\ \mathrm{ms}).
\]
快模态的 $T_2$ 由交换速率决定，而不是由小孔的 $S/V$ 决定；若按孤立孔公式换算，会得到一个并不存在的 0.16~\si{\micro m} 的“孔”。

\subsection{反算的代入示例}
\textbf{路线 A（C3）：}
\[
r_S=\frac{2\rho_2}{1/T_{2,f}-1/T_{2B}}=\frac{0.01}{1/214.79-1/3000}=2.31\ \mu\mathrm{m},\qquad r_B=\frac{0.01}{1/1086.72-1/3000}=17.04\ \mu\mathrm{m}.
\]
\textbf{路线 B（C1）：}小峰分数 $f_S=0.0495$，主峰按路线 A 得 $r_B=18.76$~\si{\micro m}：
\[
r_S=r_B\sqrt{\frac{f_S}{1-f_S}}=18.76\times\sqrt{\frac{0.0495}{0.9505}}=4.28\ \mu\mathrm{m}.
\]
\textbf{区 III 的有效半径（C5）：}主峰 $T_{2\mathrm{lm}}=1121.08$~ms，
\[
R_{\mathrm{eff}}=\frac{0.01}{1/1121.08-1/3000}=17.90\ \mu\mathrm{m}\quad\text{vs.}\quad\frac{r_B^2+r_S^2}{r_B+r_S}=\frac{400+25}{25}=17.00\ \mu\mathrm{m}.
\]
